\subsection{交错双线性型的化简。}

定理 1. --- 设 A 是一个主（交换）环，E 是有限维数为 n 的自由 A-模，\(\Phi\) 是 E 上的一个交错双线性型。则存在 E 的一个基 \((e_i)_{1 \leqslant i \leqslant n}\) 以及一个偶数 \(2r \leqslant n\) ，使得

\[
1 ^ {0} \quad \Phi (e _ {1}, e _ {2}) = \alpha_ {1}, \Phi (e _ {3}, e _ {4}) = \alpha_ {2}, \dots , \Phi (e _ {2 r - 1}, e _ {2 r}) = \alpha_ {r}
\]

其中 \(\alpha_{i}\) 是 A 中的元素 \(\neq 0\) ，且 \(\alpha_{i}\) 整除 \(\alpha_{i+1}\) （对 \(i=1,\ldots,r-1\) ）。

2° 所有其余的元素 \(\Phi(e_{i}, e_{j})\) （其中 \(i \leqslant j\) ）均为零。

理想 \(\mathrm{A}\alpha_{i}\) （ \(i=1,\ldots,r\) ）由上述条件唯一确定。E 的与 E 正交的子模 \(\mathrm{E}^{0}\) 由 \(e_{2r+1},\ldots,e_{n}\) 生成。

我们对 E 的维数 n 作归纳。n = 0 时定理是显然的。若 \(\Phi = 0\) ，定理也是显然的；因此可以假设 \(\Phi \neq 0\) 。以 f 表示线性映射 \(d_{\Phi}\) ，它从 E 映到 \(\mathbf{E}^*\) 中且在右边与 \(\Phi\) 相伴（ \(\S 1, \text{no. 1}\) ）；于是 \(f(\mathbf{E})\) 是模 \(\mathbf{E}^*\) 的一个不化为 0 的子模，后者是维数为 \(n\) 的自由模。设 \(\mathrm{A}\alpha_{1}\) 是 \(f(\mathbf{E})\) 关于 \(\mathbf{E}^*\) 的最大不变因子（第 VII 章， \(\S 4, \text{no. 2}\) ，定理 1）；我们知道（同上）存在一个基 \((e_{1}', a_{2}', \ldots, a_{n}')\) （\(\mathbf{E}^*\) 的）以及一个元素 \(f(e_{2}) \in f(\mathbf{E})\) ，使得 \(f(e_{2}) = \alpha_{1}e_{1}'\) 。设 \((e_{1}, a_{2}, \ldots, a_{n})\) 是 E 的（等同于双对偶 \(\mathbf{E}^{**}\) 的）与 \((e_{1}', a_{2}', \ldots, a_{n}')\) 对偶的基；我们有

\[
\Phi (e _ {1}, e _ {2}) = - \Phi (e _ {2}, e _ {1}) = \left\langle e _ {1}, f (e _ {2}) \right\rangle = \alpha_ {1}.\tag{1}
\]

设 P 是 E 的子模 \(Ae_{1} + Ae_{2}\) 。我们将看到，E 是 \(Ae_{1}\) 、\(Ae_{2}\) 以及与 P 正交的子模 \(P^{0}\) 的直和。为此只需证明：对每个 \(x \in E\) ，存在 A 中唯一确定的元素 \(\xi_{1}\) 、\(\xi_{2}\) ，使得 \(x - \xi_{1}e_{1} - \xi_{2}e_{2} \in P^{0}\) ，即使得

\[
\Phi (e _ {1}, x - \xi_ {1} e _ {1} - \xi_ {2} e _ {2}) = 0, \quad \Phi (e _ {2}, x - \xi_ {1} e _ {1} - \xi_ {2} e _ {2}) = 0.
\]

根据 (1)，这些条件写成

\[
\left\langle e _ {1}, f (x) \right\rangle = \xi_ {2} \alpha_ {1}, \left\langle e _ {2}, f (x) \right\rangle = - \xi_ {1} \alpha_ {1}.
\]

但是我们知道（同上），\(f(E)\) 在 \(E^{*}\) 上任一线性形式下的像都含于理想 \(A\alpha_{1}\) ，换言之，所有的值 \(\Phi(x,y)=\langle x,f(y)\rangle\) 都属于 \(A\alpha_{1}\) ；由此得到 \(\xi_{1}\) 与 \(\xi_{2}\) 的存在性与唯一性。于是 \(P^{0}\) 是秩为 \(n-2\) 的自由模；因而依归纳假设，在 \(P^{0}\) 中存在满足陈述中条件的基 \((e_{3},e_{4},\ldots,e_{n})\) 。要证明如此得到的 E 的基 \((e_{1},\ldots,e_{n})\) 也满足这些条件，只需证明 \(\alpha_{1}\) 整除 \(\alpha_{2}\) ；而这由所有的值 \(\Phi(x,y)\) 都是 \(\alpha_{1}\) 的倍式这一事实得出。于是显然 \(e_{2r+1},\ldots,e_{n}\) 生成 \(E^{0}\) 。最后，若 \((e_{i}^{\prime})\) 是 \((e_{i})\) 的对偶基，则有 \(f(e_{2j-1})=-\alpha_{j}e_{2j}^{\prime}\) 与 \(f(e_{2j})=\alpha_{j}e_{2j}^{\prime}\) （对 \(j=1,\ldots,r\) ），以及 \(f(e_{k})=0\) （对 \(k=2r+1,\ldots,n\) ）；因此诸理想 \(A\alpha_{1},A\alpha_{1},A\alpha_{2},A\alpha_{2},\ldots,A\alpha_{r},A\alpha_{r}\) 是 \(f(E)\) 关于 \(E^{*}\) 的不变因子，这就证明了它们的唯一性（第 VII 章，§ 4，第 2 小节，定理 1）。

推论 1. --- 设 A 是一个域，E 是 A 上有限维数为 n 的向量空间，Φ 是 E 上的一个交错双线性型。则存在 E 的一个基 \((e_i)_{1\leqslant i\leqslant n}\) 以及一个偶数 \(2r\leqslant n\) ，使得

\[
\Phi (\sum_ {i = 1} ^ {n} \xi_ {i} e _ {i}, \sum_ {i = 1} ^ {n} \eta_ {i} e _ {i}) = \sum_ {j = 1} ^ {r} (\xi_ {2 j - 1} \eta_ {2 j} - \xi_ {2 j} \eta_ {2 j - 1}).\tag{2}
\]

特别地，\(\Phi\) 的秩是偶数 \(2r\) 。

满足 (2) 的基称为 \(\Phi\) 的辛基。

注. --- 注意，这个推论也是 § 4，第 2 小节的命题 3 及其推论 1 的直接推论，因为按该命题的记号，由于 Φ 是交错的，必然有 H = \{0\}。

推论 2. --- 设 A 是一个域，E 是 A 上有限维数为 n 的向量空间。对每个双向量 \(z \in \bigwedge^{2} \mathrm{E}\) ，存在 E 的一个基 \((e_i)_{1 \leqslant i \leqslant n}\) ，使得

\[
z = e _ {1} \wedge e _ {2} + e _ {3} \wedge e _ {4} + \dots + e _ {2 r - 1} \wedge e _ {2 r} \quad (2 r \leqslant n).
\]

实际上，只需注意 z 典范地等同于 E* 上的一个交错双线性型（第 III 章，§ 8，第 2 小节），并对这个型应用推论 1。

把推论 1 翻译成矩阵语言，我们得到：

推论 3. --- 设 A 是一个域，R 是 A 上的一个交错方阵。R 的秩是偶数 2r，且存在 A 上的一个可逆矩阵 P，使得

\[
{ } ^ { t } P . R . P = \left( \begin{array} { c c c } 0 & I _ { r } & 0 \\ - I _ { r } & 0 & 0 \\ 0 & 0 & 0 \end{array} \right) .
\]

注. --- 若 A 是任意交换环且 R 是 A 上的奇数阶 n 交错方阵，则 det R = 0。当 A 是体时，这由推论 3 得出。当 A 是特征 ≠ 2 的体时，一个直接的证明如下：由于 \(^tR = -R\) ，我们有 det \(R = \det ^t R = (-1)^n\) det R，从而 2 det \(^tR = 0\) 。这样一来，由于交错矩阵 ( \(\alpha_{ij}\) ) 的行列式是 \(\alpha_{ij}\) （其中 i \textless{} j）的整系数多项式，代数恒等式的开拓原理（第 IV 章，§ 2，第 5 小节，附注）就表明我们的断言对任意交换环 A 都成立。

